\documentclass[10pt,a4paper]{amsart}
\usepackage[margin=19mm]{geometry}
\usepackage{amsmath,amssymb,mathtools}
\usepackage[scr=rsfs,cal=euler]{mathalfa}
\usepackage{xfrac}
\usepackage[unicode,hidelinks,bookmarks=false]{hyperref}
\newcommand{\ie}{i.e.\ }
\newcommand{\cf}{cf.\ }
\newcommand{\resp}{respectively\ }
\newcommand{\Subsection}{\subsection}
\providecommand{\nobreakdash}{\nobreak\hbox{-}}
\setlength{\emergencystretch}{2em}
\multlinegap0pt
\usepackage{bm}
\usepackage{bbm}
\usepackage{dsfont}
\usepackage{xspace}
\usepackage{cancel}
\usepackage{mathtools}
\usepackage{amsthm}
\makeatletter
\@ifundefined{thm}{\newtheorem{thm}{Thm}}{}
\@ifundefined{lem}{\newtheorem{lem}[thm]{Lemma}}{}
\makeatother
\def\bR {\mathbb{R}}
\def\bS {\mathbb{S}}
\def\cH {\mathcal{H}}
\newcommand{\can}{\mathrm{can}}
\newcommand{\dist}{\operatorname{dist}}
\newcommand{\dvol}{\operatorname{dVol}} 
\newcommand{\la}{\langle}
\newcommand{\ra}{\rangle}
\newcommand{\ud}{\mathrm{\,d}}
\newcommand{\veps}{\varepsilon}
\newcommand{\vol}{\operatorname{vol}}
\newcommand{\wto}{\rightharpoonup}
\title[Spectral Obstructions to Contracting Transport Maps on Curve]{Spectral Obstructions to Contracting Transport Maps on Curved Spaces}
\author{Shrey Aryan}
\date{Source: arXiv:2605.24705v2 (CC BY 4.0)}
\begin{document}
\maketitle

\noindent\emph{Source and licence.} Spectral Obstructions to Contracting Transport Maps on Curved Spaces. Version arXiv:2605.24705v2 (CC BY 4.0), distributed under the Creative Commons Attribution 4.0 International licence, as recorded in the arXiv licence metadata for this exact version and shown on the abstract page https://arxiv.org/abs/2605.24705v2. Full rights notice: licenses/arxiv-2605.24705-CC-BY-4.0.txt.\par\smallskip

\noindent\textit{Editorial scope.} Counted unit: the Lem whose source label is \texttt{lem:hemisphere-pancake-approximation} in the article (source0.tex lines 2519--2539, source label \texttt{lem:hemisphere-pancake-approximation}), delivered together with the article's own complete original proof of it (source0.tex lines 2541--3025, one \texttt{proof} environment of 485 lines). No mathematics is written, repaired or re-derived: statement and proof are the article's own text line for line. Required context retained uncounted: none. Every definition, display and label of those lines is retained unchanged. Omitted from this package and not redistributed: 1--2518, 2540--2540, 3026--5924. Nothing inside a retained excerpt is omitted or altered: every retained body is delivered exactly as the article's lines are written, so no omission receipt of any kind accompanies this package. Citations: the retained excerpts carry no citation command at all, so no bibliography entry of the article is transcribed and no reference is redistributed. Reference adaptation: none was needed and none was made; every reference inside the delivered unit resolves inside it, so no unresolved reference or citation is produced. Printed slips kept literal: no printed word, symbol, hypothesis or number of the counted statement or of its proof was altered. Layout and class: the article's own class is \texttt{amsart} with packages including amssymb, mathrsfs, color, mathtools, empheq, verbatim, epstopdf, pinlabel, hyperref, enumerate, bbm, graphicx, tensor, slashed; the wrapper is the lane's pinned \texttt{amsart} editorial wrapper at 10pt on A4, which restates the theorem-like environments the retained text needs and adds the article's own macro definitions that the retained text needs (\texttt{\textbackslash bR}, \texttt{\textbackslash bS}, \texttt{\textbackslash cH}, \texttt{\textbackslash can}, \texttt{\textbackslash dist}, \texttt{\textbackslash dvol}, \texttt{\textbackslash la}, \texttt{\textbackslash ra}, \texttt{\textbackslash ud}, \texttt{\textbackslash veps}, \texttt{\textbackslash vol}, \texttt{\textbackslash wto}), with extra wrapper packages (bm, bbm, dsfont, xspace, cancel, mathtools). The delivered numbering is the unit's own. Assets: no retained span contains a figure environment, a graphics-inclusion command, a PDF-page inclusion, a table or a TikZ picture; no image file is copied into this package, the package declares no asset dependency and no image file is redistributed with it. Licence: version arXiv:2605.24705v2 of this article is distributed under the Creative Commons Attribution 4.0 International licence, as recorded in the arXiv licence metadata for this exact version and shown on the abstract page https://arxiv.org/abs/2605.24705v2. A full rights notice is delivered at licenses/arxiv-2605.24705-CC-BY-4.0.txt. Characters: every retained excerpt is pure ASCII, so the delivered engine, XeLaTeX with its default Latin Modern text font, reports no missing character.
\begin{lem}\label{lem:hemisphere-pancake-approximation}
Let $d\geq2$, let $\bS^d_+$ be a closed hemisphere of the equatorial copy
\begin{align}
        \bS^d
        =
        \bS^{d+1}\cap\{z=0\},
\end{align}
and let $K\subset\bS^d_+$ be a closed geodesically convex set satisfying $ \vol_{g_{\can}}(K)>0.$ Then there exist smooth closed strictly convex hypersurfaces $\Sigma_j \subset (\bS^{d+1},\overline g_{\can})$ contained in an open hemisphere and invariant under $R(x,z)=(x,-z)$ such that, setting $ \Sigma_j^+:=\Sigma_j\cap\{z\geq0\},$ we have
\begin{align}\label{eq:hemisphere-pancake-Hausdorff-limit}
        d_H(\Sigma_j^+,K)
        \rightarrow
        0,
\end{align}
and
\begin{align}\label{eq:hemisphere-pancake-measure-limit}
        \dvol_{g_{\Sigma_j}}|_{\Sigma_j^+}
        \wto
        \mathbf{1}_K\dvol_{g_{\can}}
\end{align}
as $j\to \infty$. Moreover, $ \Sigma_j\cap\{z>0\}$ is connected for every $j$. Here $d_H$ denotes the Hausdorff distance between compact subsets of $(\bS^{d+1},\overline g_{\can})$.
\end{lem}

\begin{proof}
We first construct an approximation of $K$ that is away from the boundary of the hemisphere. Let $N\in\bS^d$ denote the pole of the hemisphere $\bS^d_+$, so that
\begin{align}
        \bS^d_+
        =
        \{x\in\bS^d:x\cdot N\geq0\}.
\end{align}
Since $K$ is geodesically convex and has positive $d$-dimensional volume,
it has nonempty interior. Choose $p\in\operatorname{int}(K)$. Since $K\subset\bS^d_+$, we have
$p\cdot N>0$. For each $j\geq1$, set
\begin{align}
        \delta_j
        :=
        \frac{p\cdot N}{j+1},
        \qquad
        K_j
        :=
        K\cap\{x\cdot N\geq\delta_j\}.
\end{align}
Then the sets $K_j$ are compact, increasing, geodesically convex, and
contained in the open hemisphere $\{x\cdot N>0\}$. Moreover,
$p\in\operatorname{int}(K_j)$ for every $j\geq1$. We claim that
\begin{align}
        \overline{\bigcup_{j=1}^\infty K_j}
        =
        K.
\end{align}
Indeed, every $q\in K$ satisfying $q\cdot N>0$ belongs to $K_j$ for all
sufficiently large $j$. If instead $q\cdot N=0$, then $p$ and $q$ are not antipodal since $p\cdot N>0$. Thus, if we define, for $0\leq s<1$,
\begin{align}
        q_s
        :=
        \frac{(1-s)p+sq}{|(1-s)p+sq|},
\end{align}
then the point $q_s$ lies on the minimizing geodesic segment from $p$ to $q$ and
therefore belongs to $K$ by geodesic convexity. Furthermore,
\begin{align}
        q_s\cdot N
        =
        \frac{(1-s)p\cdot N}{|(1-s)p+sq|}
        >
        0,
\end{align}
so $q_s\in K_j$ for all sufficiently large $j$. Since $q_s\to q$ as
$s\to1$, the claim follows. Since $K_j\subset K$, the preceding claim and the compactness of $K$ imply
\begin{align}
        d_H(K_j,K)
        \rightarrow
        0
\end{align}
as $j\to \infty.$ Moreover, as $j\to \infty$
\begin{align}
        \mathbf{1}_{K_j}
        \rightarrow
        \mathbf{1}_K
        \quad
        \dvol_{g_{\can}}\text{-almost everywhere},
\end{align}
because the equator $\{x\cdot N=0\}$ has zero $d$-dimensional volume. Consequently, by the dominated convergence theorem,
\begin{align}\label{eq:hemisphere-pancake-truncation-limit}
        \mathbf{1}_{K_j}\dvol_{g_{\can}}
        \wto
        \mathbf{1}_K\dvol_{g_{\can}}.
\end{align}
Next, we further approximate the sets $K_j$ by convex sets with smooth uniformly
convex boundaries. Regard $N$ also as the point $(N,0)\in\bS^{d+1}\subset\bR^{d+1}\times\bR$,
and set $\cH_N:=\{Y\in\bS^{d+1}:Y\cdot N>0\}.$ Identify
\begin{align}
        T_N\bS^{d+1}
        =
        T_N\bS^d\oplus\operatorname{span}\{\partial_z\}
        \simeq
        \bR^d\times\bR.
\end{align}
Define the gnomonic projection based at $N$ by
\begin{align}
        \Gamma:\cH_N
        \rightarrow
        T_N\bS^{d+1},\quad 
        \Gamma(Y)
        :=
        \frac{Y}{Y\cdot N}-N.
\end{align}
Its inverse is
\begin{align}
        G:T_N\bS^{d+1}
        \rightarrow
        \cH_N,\quad 
        G(V)
        :=
        \frac{N+V}{\sqrt{1+|V|^2}}.
\end{align}
The restriction of $\Gamma$ to
$\cH_N\cap\{z=0\}$ takes values in
$T_N\bS^d\simeq\bR^d$. Since
\begin{align}
        K_j
        \subset
        \{x\in\bS^d:x\cdot N\geq\delta_j\}
        \subset
        \cH_N,
\end{align}
we may define
\begin{align}
        A_j
        :=
        \Gamma(K_j)
        \subset
        T_N\bS^d
        \simeq
        \bR^d.
\end{align}
Because the gnomonic projection maps geodesic segments contained in
$\cH_N$ onto Euclidean line segments, $A_j$ is compact and convex.
Furthermore, since $\Gamma$ is a diffeomorphism and $K_j$ has nonempty
interior, $A_j$ has nonempty interior. We next smooth $A_j$. Let
\begin{align}
        h_j(u)
        :=
        \sup_{x\in A_j}x\cdot u,
        \quad 
        f_j(x)
        :=
        \max_{u\in\bS^{d-1}}\bigl(x\cdot u-h_j(u)\bigr).
\end{align}
Then $f_j$ is convex and $1$-Lipschitz,
\begin{align}
        A_j
        =
        \{f_j\leq0\},
        \qquad
        \max\{f_j,0\}
        =
        \dist(\cdot,A_j).
\end{align}
Fix a nonnegative standard mollifier $\varphi_\eta$ supported in
$B_\eta(0)$ and satisfying $ \int_{\bR^d}\varphi_\eta\,\ud x = 1$. Choose $R_j$ such that $A_j\subset B_{R_j}(0)$, and define
\begin{align}
        F_{j,\eta}
        :=(f_j*\varphi_\eta)+\eta|x|^2,
        \quad 
        c_{j,\eta}
        :=2\eta(1+R_j^2),
        \\
        A_{j,\eta}
        :=\{F_{j,\eta}\leq c_{j,\eta}\},
        \quad 
        \rho_{j,\eta}
        :=F_{j,\eta}-c_{j,\eta}.
\end{align}
Since $\|f_j*\varphi_\eta-f_j\|_{C^0}\leq\eta$, we have
\begin{align}\label{eq:hemisphere-pancake-smooth-approximation}
        A_j
        \subset
        A_{j,\eta}
        \subset
        A_j+(c_{j,\eta}+\eta)\overline B_1(0),
        \qquad
        D^2\rho_{j,\eta}
        \geq
        2\eta I.
\end{align}
Since $f_j<0$ at every interior point of $A_j$, for all sufficiently small
$\eta>0$ the level $c_{j,\eta}$ lies strictly above the unique minimum of
the strictly convex function $F_{j,\eta}$. Hence $\nabla\rho_{j,\eta}\neq0$ on $\partial A_{j,\eta}$. Therefore, $A_{j,\eta}$ is a smooth uniformly convex body and
$d_H(A_{j,\eta},A_j)\to0$ as $\eta\to0$. Set
\begin{align}
        S_{j,\eta}
        :=
        G(A_{j,\eta}\times\{0\}).
\end{align}
Let $G_0(v)=G(v,0)$ for $v\in T_N \bS^{d}.$ For $v,\xi\in T_N\bS^d$, we have
\begin{align}
        d(G_0)_v(\xi)
        =
        \frac{(\xi,0)}{\sqrt{1+|v|^2}}
        -
        \frac{\la v,\xi\ra(N+v,0)}
        {(1+|v|^2)^{3/2}}.
\end{align}
It follows that
\begin{align}
        (G_0^*g_{\can})_v(\xi,\zeta)
        =
        \frac{\la\xi,\zeta\ra}{1+|v|^2}
        -
        \frac{\la v,\xi\ra\la v,\zeta\ra}
        {(1+|v|^2)^2}.
\end{align}
Therefore,
\begin{align}\label{eq:hemisphere-pancake-gnomonic-Jacobian}
        G_0^*\dvol_{g_{\can}}
        =
        J_{G_0}(v)\,\ud v,
        \qquad
        J_{G_0}(v)
        :=
        (1+|v|^2)^{-\frac{d+1}{2}}.
\end{align}

Fix $j$. Since
$d_H(A_{j,\eta},A_j)\to0$, all the sets $A_{j,\eta}$ with $\eta>0$
sufficiently small are contained in a fixed compact subset of
$T_N\bS^d$. The map $G_0$ is Lipschitz on this compact set, and hence
\begin{align}
        d_H(S_{j,\eta},K_j)
        \leq
        C_jd_H(A_{j,\eta},A_j)
        \rightarrow
        0,
\end{align}
as $\eta\to 0$. Moreover, by
\eqref{eq:hemisphere-pancake-smooth-approximation}, we have
\begin{align}
        A_j
        \subset
        A_{j,\eta}
        \subset
        A_j+r_{j,\eta}\overline B_1(0),
        \qquad
        r_{j,\eta}
        :=
        c_{j,\eta}+\eta
        \rightarrow
        0,
\end{align}
as $\eta\to 0.$ Consequently, as $\eta \to 0$
\begin{align}
        \mathbf{1}_{A_{j,\eta}}(v)
        \rightarrow
        \mathbf{1}_{A_j}(v)
        \quad
        \text{for every }v\in T_N\bS^d.
\end{align}
Let $f\in C^0(\bS^{d+1})$. Using
\eqref{eq:hemisphere-pancake-gnomonic-Jacobian}, we obtain
\begin{align}
        \int_{S_{j,\eta}}f\,\dvol_{g_{\can}}
        =
        \int_{T_N\bS^d}
        f(G_0(v))
        \mathbf{1}_{A_{j,\eta}}(v)
        J_{G_0}(v)\,\ud v.
\end{align}
Thus, the dominated convergence theorem gives
\begin{align}
        \int_{S_{j,\eta}}f\,\dvol_{g_{\can}}
        \rightarrow
        \int_{T_N\bS^d}
        f(G_0(v))
        \mathbf{1}_{A_j}(v)
        J_{G_0}(v)\,\ud v
        =
        \int_{K_j}f\,\dvol_{g_{\can}}.
\end{align}
Therefore, $ \mathbf{1}_{S_{j,\eta}}\dvol_{g_{\can}} \wto \mathbf{1}_{K_j}\dvol_{g_{\can}}$ as $\eta \to 0$. Choose $0<\eta_j<1/j$ diagonally so that, setting
\begin{align}
        A'_j
        :=
        A_{j,\eta_j},
        \quad 
        \rho_j
        :=
        \rho_{j,\eta_j},
        \quad 
        S_j
        :=
        S_{j,\eta_j},
\end{align}
we have as $j\to \infty$
\begin{align}\label{eq:hemisphere-pancake-smooth-limit}
        d_H(S_j,K)
        \rightarrow
        0,
        \quad
        \mathbf{1}_{S_j}\dvol_{g_{\can}}
        \wto
        \mathbf{1}_K\dvol_{g_{\can}}.
\end{align}
We now thicken each $S_j$ symmetrically in the additional $z$-direction
to obtain a smooth closed strictly convex hypersurface in $\bS^{d+1}$. Here the reader should think about the simple case of the unit interval $[-1,1]$ which can be fattened to obtain an ellipse that collapses onto the plane from upper and lower side. Thus, for $\varepsilon>0$, define
\begin{align}
        B_{j,\varepsilon}
        :=
        \left\{
        (x,z)\in\bR^d\times\bR:
        \rho_j(x)+\frac{z^2}{\varepsilon^2}\leq0
        \right\}.
\end{align}
The Hessian of the function $\rho_j+z^2/\veps^2$ is
\begin{align}
        \begin{pmatrix}
        D^2\rho_j&0\\
        0&2\varepsilon^{-2}
        \end{pmatrix}
        >
        0,
\end{align}
and its gradient does not vanish on the zero level. Thus
$\widehat\Sigma_{j,\varepsilon}:=\partial B_{j,\varepsilon}$ is a smooth
closed strictly convex Euclidean hypersurface invariant under
$(x,z)\mapsto(x,-z)$. Now observe that if $\zeta$ is the Euclidean unit normal to
$\widehat\Sigma_{j,\varepsilon}$ at $y$, then for tangent vectors $v,w$,
\begin{align}
        \operatorname{II}_{\bS^{d+1}}(dG_y(v),dG_y(w))
        =
        \frac{
        \operatorname{II}_{\bR^{d+1}}(v,w)
        }{
        \sqrt{1+|y|^2}
        \sqrt{1+(y\cdot\zeta)^2}
        }.
\end{align}
Thus, the map $G$ preserves positivity of the second fundamental form and therefore, $\Sigma_{j,\varepsilon} := G(\widehat\Sigma_{j,\varepsilon})$ is a smooth closed strictly convex hypersurface in an open hemisphere and is
invariant under $R$.

The upper half of $\widehat\Sigma_{j,\varepsilon}$ is the graph of $ z_{j,\varepsilon}(x)= \varepsilon\sqrt{-\rho_j(x)}$ for $x\in A_j'$ and hence $\Sigma_{j,\varepsilon}^+$ is parametrized by
\begin{align}
        \Psi_{j,\varepsilon}:\operatorname{int}(A'_j)
        \rightarrow
        \Sigma_{j,\varepsilon}\cap\{z>0\},
        \quad 
        \Psi_{j,\varepsilon}(x)
        :=
        G\bigl(x,z_{j,\varepsilon}(x)\bigr).
\end{align}
The same formula extends $\Psi_{j,\varepsilon}$ continuously to $A'_j$;
we continue to denote this extension by $\Psi_{j,\varepsilon}$. Since
$G$ preserves the sign of the last coordinate, we have
\begin{align}
        \Psi_{j,\varepsilon}
        \left(
        \operatorname{int}(A'_j)
        \right)
        =
        \Sigma_{j,\varepsilon}\cap\{z>0\},
        \quad
        \Psi_{j,\varepsilon}(A'_j)
        =
        \Sigma_{j,\varepsilon}^+.
\end{align}
 We have
\begin{align}
        \widehat\Sigma_{j,\varepsilon}\cap\{z>0\}
        =
        \{
        (x,\varepsilon\sqrt{-\rho_j(x)}):
        x\in\operatorname{int}(A'_j)
        \}.
\end{align}
Since $A'_j$ is convex, its interior is connected. The inverse gnomonic
projection $G$ preserves the sign of the $z$-coordinate, and hence
$\Sigma_{j,\varepsilon}\cap\{z>0\}$ is connected. Since $\rho_j=0$ and $\nabla\rho_j\neq0$ on $\partial A'_j$, there exist
constants $c_j,C_j,r_j>0$ such that
\begin{align}\label{eq:hemisphere-pancake-defining-function-distance}
        c_j\dist(x,\partial A'_j)
        \leq
        -\rho_j(x)
        \leq
        C_j\dist(x,\partial A'_j)
\end{align}
whenever $x\in A'_j$ and
$\dist(x,\partial A'_j)<r_j$. Moreover,
for $x\in\operatorname{int}(A'_j)$,
\begin{align}
        \nabla z_{j,\varepsilon}(x)
        =
        -\frac{\varepsilon\nabla\rho_j(x)}
        {2\sqrt{-\rho_j(x)}}.
\end{align}
Consequently, after increasing $C_j$ if necessary,
\begin{align}\label{eq:hemisphere-pancake-graph-gradient-bound}
        |\nabla z_{j,\varepsilon}(x)|
        \leq
        \frac{C_j\varepsilon}
        {\sqrt{\dist(x,\partial A'_j)}}
        \leq
        \frac{C_j}
        {\sqrt{\dist(x,\partial A'_j)}}
\end{align}
for $0<\varepsilon\leq1$. In particular, for every fixed
$x\in\operatorname{int}(A'_j)$, $z_{j,\varepsilon}(x)
        \rightarrow
        0,
        |\nabla z_{j,\varepsilon}(x)|
        \rightarrow
        0$ as $\varepsilon\to0$. Since $A'_j$ is compact, $\sup_{x\in A'_j}|z_{j,\varepsilon}(x)| \leq C_j\varepsilon.$ The map $G$ is Lipschitz on a fixed compact set containing all the graphs
for $0<\varepsilon\leq1$. Recalling that $ S_j
        =
        G(A'_j\times\{0\}),$ we obtain
\begin{align}\label{eq:hemisphere-pancake-upper-sheet-Hausdorff}
        d_H(\Sigma_{j,\varepsilon}^+,S_j)
        \leq
        \sup_{x\in A'_j}
        d_{\overline g_{\can}}
        \left(
        G\bigl(x,z_{j,\varepsilon}(x)\bigr),
        G(x,0)
        \right)
        \leq
        C_j\varepsilon
        \rightarrow
        0,
\end{align}
as $\varepsilon\to 0.$ We next prove convergence of the volume measures. First note that, for $y\in T_N\bS^{d+1}$ and $v,w\in T_N\bS^{d+1}$
\begin{align}
        (G^*\overline g_{\can})_y(v,w)
        =
        \frac{\la v,w\ra}{1+|y|^2}
        -
        \frac{\la y,v\ra\la y,w\ra}
        {(1+|y|^2)^2}.
\end{align}
Applying this formula with
\begin{align}
        y
        =
        \bigl(x,z_{j,\varepsilon}(x)\bigr),
        \qquad
        v
        =
        \bigl(\xi,\nabla z_{j,\varepsilon}(x)\cdot\xi\bigr),
\end{align}
and computing the determinant of the induced metric gives
\begin{align}\label{eq:hemisphere-pancake-spherical-graph-Jacobian}
        J_{j,\varepsilon}(x)
        :=
        \sqrt{
        \det\bigl(\Psi_{j,\varepsilon}^*
        \overline g_{\can}\bigr)
        }
        =
        \frac{
        \sqrt{
        1+|\nabla z_{j,\varepsilon}(x)|^2
        +
        \left(
        z_{j,\varepsilon}(x)
        -
        x\cdot\nabla z_{j,\varepsilon}(x)
        \right)^2
        }
        }{
        \left(
        1+|x|^2+z_{j,\varepsilon}(x)^2
        \right)^{\frac{d+1}{2}}
        }.
\end{align}
Thus, for every $x\in\operatorname{int}(A'_j)$, $ J_{j,\varepsilon}(x) \to \frac{1}{(1+|x|^2)^{\frac{d+1}{2}}}$ as $\varepsilon \to 0.$ Since $A'_j$ is bounded, it follows from \eqref{eq:hemisphere-pancake-graph-gradient-bound} that
\begin{align}
        J_{j,\varepsilon}(x)
        \leq
        C_j ( 1+
        \dist(x,\partial A'_j)^{-1/2}).
\end{align}
The right-hand side is integrable on $A'_j$.  Indeed, in a smooth collar neighborhood of the boundary we have
\begin{align}
        \int_0^{r_j}s^{-1/2}\,\ud s
        <
        \infty.
\end{align}
Since $\partial A'_j$ has zero $d$-dimensional measure, the area formula
on $\operatorname{int}(A'_j)$ gives for any $f\in C^0(\bS^{d+1})$
\begin{align}
        \int_{\Sigma_{j,\varepsilon}^+}
        f\,\dvol_{g_{\Sigma_{j,\varepsilon}}}
        =
        \int_{A'_j}
        f\bigl(\Psi_{j,\varepsilon}(x)\bigr)
        J_{j,\varepsilon}(x)\,\ud x.
\end{align}
Thus the dominated convergence theorem implies
\begin{align}
        \int_{\Sigma_{j,\varepsilon}^+}
        f\,\dvol_{g_{\Sigma_{j,\varepsilon}}}
        \rightarrow
        \int_{A'_j}
        f(G(x,0))
        \frac{\ud x}{(1+|x|^2)^{\frac{d+1}{2}}}
        =
        \int_{S_j}f\,\dvol_{g_{\can}}.
\end{align}
Therefore, $ \dvol_{g_{\Sigma_{j,\varepsilon}}} |_{\Sigma_{j,\varepsilon}^+} \wto \mathbf{1}_{S_j}\dvol_{g_{\can}}$ as $\veps\to 0.$ Choosing $0<\varepsilon_j<1/j$ diagonally and setting $\Sigma_j:=\Sigma_{j,\varepsilon_j}$, the conclusion follows from
\eqref{eq:hemisphere-pancake-smooth-limit}.
\end{proof}

\end{document}
